
\documentclass[11pt]{article}

\usepackage[T1]{fontenc}
\usepackage{lmodern}
\usepackage[margin=1.05in]{geometry}
\usepackage{amsmath,amssymb,amsthm,mathtools}
\usepackage{enumitem}
\usepackage{booktabs}
\usepackage{microtype}
\usepackage[colorlinks=true,linkcolor=blue,citecolor=blue,urlcolor=blue]{hyperref}
\hypersetup{
  pdftitle={Erdos--Turan, Berry--Esseen, and Edgeworth Universality after One Shuffle},
  pdfauthor={Christopher D. Long},
  pdfsubject={Random permutations, riffle shuffles, shelf shuffles, and permutation order}
}

\newtheorem{theorem}{Theorem}[section]
\newtheorem{proposition}[theorem]{Proposition}
\newtheorem{lemma}[theorem]{Lemma}
\newtheorem{corollary}[theorem]{Corollary}
\newtheorem{inputtheorem}[theorem]{Input theorem}
\newtheorem{criterion}[theorem]{Criterion}
\theoremstyle{definition}
\newtheorem{definition}[theorem]{Definition}
\theoremstyle{remark}
\newtheorem{remark}[theorem]{Remark}

\newcommand{\Pp}{\mathbb P}
\newcommand{\Ee}{\mathbb E}
\newcommand{\Var}{\operatorname{Var}}
\newcommand{\Cov}{\operatorname{Cov}}
\newcommand{\one}{\mathbf 1}
\newcommand{\dd}{\,\mathrm d}
\newcommand{\ord}{\operatorname{ord}}
\newcommand{\lcm}{\operatorname{lcm}}
\newcommand{\Cyc}{\operatorname{Cyc}}
\newcommand{\Pois}{\operatorname{Pois}}
\newcommand{\Geo}{\operatorname{Geo}_0}
\newcommand{\TV}{\mathrm{TV}}
\newcommand{\Kol}{\mathrm d_{\mathrm K}}
\newcommand{\cL}{\mathcal L}
\newcommand{\cO}{\mathcal O}
\newcommand{\cZ}{\mathcal Z}
\newcommand{\cQ}{\mathcal Q}
\newcommand{\cR}{\mathcal R}
\newcommand{\cG}{\mathcal G}
\newcommand{\R}{\mathbb R}
\newcommand{\N}{\mathbb N}
\newcommand{\eps}{\varepsilon}

\title{Erd\H{o}s--Tur\'an, Berry--Esseen, and Edgeworth\\
Universality after One Shuffle}
\author{Christopher D. Long}
\date{Edgeworth and optimality working branch}

\begin{document}
\maketitle

\begin{abstract}
We prove functional Erd\H{o}s--Tur\'an universality, a sharp
Berry--Esseen theorem, and a universal first Edgeworth expansion for
several nonexchangeable random-permutation models generated by a single
shuffle.  The common input is an exponentially Ewens-stable cycle index:
under geometric randomization of the size, its logarithm differs from the
uniform-permutation cycle index by coefficients that decay exponentially
in the cycle length.  For every such triangular family, if $C_{n,r}$
denotes the number of $r$-cycles, then
\[
 \begin{aligned}
 &\left(
 \frac{\sum_{r\le n^t}C_{n,r}-t\log n}{\sqrt{\log n}},
 \frac{\log\lcm\{r\le n^t:C_{n,r}>0\}
             -\frac12t^2(\log n)^2}
      {\sqrt{\frac13(\log n)^3}}
 \right)_{0\le t\le1}
 \\
 &\hspace{35mm}\Rightarrow
 \left(
 W(t),\sqrt3\int_0^t u\,\dd W(u)
 \right)_{0\le t\le1}
 \end{aligned}
\]
in $D([0,1],\R^2)$.  In particular, the endpoint correlation between the
number of cycles and the logarithm of the order is $\sqrt3/2$.

Writing $Y_n=\log\ord(\pi_n)$, $m_n=\Ee Y_n$, and $\ell_n=\log n$, we
prove, uniformly over every subfamily with common stability constants,
\[
 \sup_{x\in\R}\left|
 \Pp\left\{\frac{Y_n-m_n}{\sqrt{\ell_n^3/3}}\le x\right\}
 -\Phi(x)
 -\frac{\sqrt3}{8\sqrt{\ell_n}}(1-x^2)\varphi(x)
 \right|
 \le
 C\left(\frac{\log\ell_n}{\ell_n}\right)^{2/3}.
\]
Consequently,
\[
 \Kol\!\left(
 \frac{Y_n-m_n}{\sqrt{\ell_n^3/3}},N(0,1)
 \right)
 =
 \frac{\sqrt3}{8\sqrt{2\pi\ell_n}}
 +O\!\left(\left(\frac{\log\ell_n}{\ell_n}\right)^{2/3}\right),
\]
so the $(\log n)^{-1/2}$ rate is genuinely optimal, with a universal
leading constant.  We also retain the explicit Barbour--Tavar\'e form
centered at $\frac12\ell_n^2-\ell_n\log\ell_n$.

The proof is a quantitative long-cycle universality argument.  Exponential
Ewens stability implies that, after deleting cycles of length at most $b$,
the entire remaining cycle vector is within $O(\rho^b)$ in total variation
of the corresponding vector for a uniform permutation.  A cutoff
$b\asymp\log\log n$ transfers the sharp Berry--Esseen upper bound.  For the
Edgeworth expansion we take $b\asymp\log n$, making the coupling failure
polynomially small in $n$; the deleted short cycles contribute only
$O((\log\log n)^2)$ in mean.  Zacharovas's first Edgeworth expansion for
the uniform permutation then transfers with an error
$o((\log n)^{-1/2})$.

We verify the criterion for random standardized permutations generated by
i.i.d.\ letters with maximal atom uniformly bounded away from one.  This
gives all results for arbitrary biased riffle shuffles under the same
nonconcentration condition.  In particular, a single ordinary $2$-riffle
shuffle already has the classical functional law, the sharp
Berry--Esseen rate, and the universal first Edgeworth correction, although
its total-variation distance from uniform tends to one.  The same
conclusions hold for major-index-biased permutations in the stated uniform
parameter range.  Using the cycle-index formula of Tripathi, we also prove
the results for one asymmetric shelf shuffle, uniformly over an arbitrary
number of shelves and over asymmetry parameters bounded away from zero and
one.

The arithmetic step is self-contained.  We extract a weighted
additive-to-least-common-multiple transfer theorem from the von Mangoldt
decomposition.  Harmonic one-cycle bounds and harmonic two-cycle factorial
bounds suffice for the functional law, while the quantitative scalar
results follow from the stronger long-cycle coupling.
\end{abstract}

\medskip
\noindent\textbf{2020 Mathematics Subject Classification.}
Primary 60C05; secondary 60F17, 05A05, 60B15.

\smallskip
\noindent\textbf{Keywords.}
random permutations; Erd\H{o}s--Tur\'an law; riffle shuffle; shelf shuffle;
major index; cycle index; functional central limit theorem; Berry--Esseen
bound; Edgeworth expansion; optimal rate; least common multiple.

\tableofcontents

\section{Introduction}

All asymptotic statements are as $n\to\infty$.  We write $\Phi$ and $\varphi$ for the standard normal distribution function and density.  For a permutation $\pi\in S_n$, its order is
\[
 \ord(\pi)=\lcm\{\text{cycle lengths of }\pi\}.
\]
Erd\H{o}s and Tur\'an proved that for a uniform random permutation
$\Pi_n$,
\[
 \frac{\log\ord(\Pi_n)-\frac12(\log n)^2}
 {\sqrt{\frac13(\log n)^3}}
 \Rightarrow N(0,1)
 \qquad\text{\cite{ErdosTuran1967}.}
\]
Barbour and Tavar\'e proved the sharp quantitative form: with
$\ell_n=\log n$,
\[
 \Kol\!\left(
 \frac{\log\ord(\Pi_n)-\frac12\ell_n^2
       +\ell_n\log\ell_n}
      {\sqrt{\ell_n^3/3}},
 N(0,1)
 \right)
 =O(\ell_n^{-1/2});
\]
see \cite{BarbourTavare1994} and the explicit formulation recorded in
\cite{Zacharovas2009}.  Zacharovas subsequently proved the first uniform
Edgeworth expansion, centered by the exact mean, with correction
\[
 \frac{\sqrt3}{8\sqrt{\log n}}(1-x^2)\varphi(x)
\]
and remainder $O(((\log\log n)/(\log n))^{2/3})$; in particular, the
$(\log n)^{-1/2}$ rate is optimal.  The subject was placed in the
probabilistic theory of logarithmic combinatorial structures; see, among
others, \cite{ArratiaTavare1992,ABT2003}.  The classical proof reflects two
facts.  First, cycle counts on logarithmic scales behave like independent
Poisson variables of means $1/r$.  Second, the logarithm of the product of
the cycle lengths can be replaced by the logarithm of their least common
multiple.

The present paper shows that this mechanism can be present long before the
underlying permutation is globally mixed.  Our most concrete result is the
following.

\begin{theorem}[One ordinary $2$-riffle shuffle]\label{thm:headline}
Let $\rho_n$ be the permutation obtained from the identity by one ordinary
Gilbert--Shannon--Reeds $2$-riffle shuffle.  Put
\[
 \ell_n=\log n,\qquad
 Y_n=\log\ord(\rho_n),\qquad
 m_n=\Ee Y_n,\qquad
 s_n=\sqrt{\ell_n^3/3}.
\]
Then the bivariate functional theorem in
Theorem~\ref{thm:general-functional} holds, and
\[
 \Kol\!\left(
 \frac{Y_n-\frac12\ell_n^2+\ell_n\log\ell_n}{s_n},
 N(0,1)
 \right)
 \le \frac{C}{\sqrt{\ell_n}}.
\]
Moreover,
\[
 \sup_{x\in\R}\left|
 \Pp\left\{\frac{Y_n-m_n}{s_n}\le x\right\}
 -\Phi(x)
 -\frac{\sqrt3}{8\sqrt{\ell_n}}(1-x^2)\varphi(x)
 \right|
 \le
 C\left(\frac{\log\ell_n}{\ell_n}\right)^{2/3},
\]
and therefore
\[
 \Kol\!\left(\frac{Y_n-m_n}{s_n},N(0,1)\right)
 =\frac{\sqrt3}{8\sqrt{2\pi\ell_n}}
 +O\!\left(\left(\frac{\log\ell_n}{\ell_n}\right)^{2/3}\right).
\]
Thus the normal-approximation rate is optimal.  Nevertheless,
\[
 d_{\TV}\bigl(\cL(\rho_n),\operatorname{Unif}(S_n)\bigr)\longrightarrow1.
\]
\end{theorem}

Thus the order statistic has already reached its uniform-permutation
universality class after one minimally nontrivial riffle shuffle, even
though the permutation itself remains asymptotically maximally far from
uniform.

The proof is organized around a cycle-index criterion that applies to two
different mechanisms.

\begin{enumerate}[label=\textnormal{(\roman*)},leftmargin=2.2em]
\item For standardized permutations, Guerder's primitive-word theorem
gives an exact geometric-size decomposition of the cycle types
\cite{Guerder2026}.  It yields a cycle index whose non-Ewens coefficients
decay exponentially whenever the largest letter probability is bounded
away from one.

\item For asymmetric shelf shuffles, Tripathi's exact cycle-index formula
has the same analytic form \cite{Tripathi2026}.  The unique nondecaying
coefficient is precisely the uniform-permutation coefficient; every other
coefficient decays exponentially when the asymmetry parameter stays away
from zero and one.
\end{enumerate}

The common probabilistic theorem is stronger than an endpoint central
limit theorem.  Put
\[
 K_n(t)=\sum_{r\le n^t}C_{n,r},
 \qquad
 \cO_n(t)=\log\lcm\{r\le n^t:C_{n,r}>0\},
 \qquad 0\le t\le1.
\]
We prove
\[
 \left(
 \frac{K_n(t)-t\log n}{\sqrt{\log n}},
 \frac{\cO_n(t)-\frac12t^2(\log n)^2}
      {\sqrt{\frac13(\log n)^3}}
 \right)
 \Rightarrow
 \left(
 W(t),\sqrt3\int_0^t u\,\dd W(u)
 \right).
\]
The two coordinates are driven by the same Brownian noise on logarithmic
cycle scale.  At the endpoint this gives
\[
 \left(
 \frac{K_n-\log n}{\sqrt{\log n}},
 \frac{\log\ord(\pi_n)-\frac12(\log n)^2}
      {\sqrt{\frac13(\log n)^3}}
 \right)
 \Rightarrow
 N\!\left(
 0,
 \begin{pmatrix}
 1&\sqrt3/2\\
 \sqrt3/2&1
 \end{pmatrix}
 \right).
\]

A further contribution is an abstract arithmetic transfer theorem.
For arbitrary component multiplicities $(C_r)$, let
\[
 B=\sum_r C_r\log r,
 \qquad
 L=\lcm\{r:C_r>0\}.
\]
Writing $D_m=\sum_{m\mid r}C_r$, the exact identity
\[
 B-\log L=\sum_m\Lambda(m)(D_m-1)_+
\]
reduces the product--LCM discrepancy to repeated prime-power divisibility.
We show that the bounds
\[
 \Ee C_r\ll\frac1r,
 \qquad
 \Ee\!\left[C_r(C_s-\one_{\{r=s\}})\right]\ll\frac1{rs}
\]
already imply
\[
 \Ee(B-\log L)=O(\log n\log\log n).
\]
The same estimate transfers the entire process, because the truncated
product--LCM discrepancy is nondecreasing in the truncation parameter.

\subsection*{Main applications}

The standardized-permutation, sharp Berry--Esseen, and Edgeworth theorems allow an arbitrary triangular array
of discrete distributions $p^{(n)}$ on finite or countable ordered
alphabets, subject only to
\[
 \sup_n\|p^{(n)}\|_\infty<1.
\]
No limiting alphabet distribution is assumed.  It contains:

\begin{itemize}[leftmargin=2em]
\item one biased riffle shuffle with arbitrary pile probabilities satisfying
the same nonconcentration condition;
\item one ordinary $a_n$-riffle shuffle for every sequence $a_n\ge2$;
\item major-index-biased permutations with every fixed parameter
$q\in(0,1)$, and uniformly for $q_n\ge q_0>0$.
\end{itemize}

The shelf-shuffle theorem permits arbitrary shelf numbers $m_n\ge1$ and
parameters
\[
 \eps\le p_n\le1-\eps.
\]
In particular, it covers a single shelf and a growing number of shelves in
one uniform statement.

\subsection*{Relation with earlier work}

The cycle theory of riffle shuffles goes back to
Diaconis--McGrath--Pitman \cite{DiaconisMcGrathPitman1995}, and generalized
riffle distributions were developed by Stanley \cite{Stanley2001}.
Guerder recently gave a common primitive-word theory for standardized
permutations, including riffle and major-index models
\cite{Guerder2026}.  The symmetric shelf shuffle was analyzed by
Diaconis--Fulman--Holmes \cite{DiaconisFulmanHolmes2013}; Tripathi extended
its cycle index to the asymmetric model \cite{Tripathi2026}.  The present
results use those exact enumerative inputs but prove the functional
Erd\H{o}s--Tur\'an laws and the arithmetic transfer independently.

The weighted additive-to-LCM argument was motivated by the comparable-rate
Luce order problem.  No theorem from the Luce papers is used here.

\section{Weighted additive-to-LCM transfer}

We begin with the model-independent arithmetic step.

Let $(C_r)_{1\le r\le U}$ be arbitrary nonnegative integer-valued random
variables.  Define
\[
 B(U)=\sum_{r\le U}C_r\log r,
 \qquad
 L(U)=\lcm\{r\le U:C_r>0\},
\]
with $\lcm\varnothing=1$, and
\[
 D_m(U)=\sum_{\substack{r\le U\\m\mid r}}C_r.
\]

\begin{theorem}[Exact weighted transfer]\label{thm:abstract-transfer}
For every $1\le y\le U$,
\[
 B(U)-\log L(U)
 =
 \sum_{m\le U}\Lambda(m)\bigl(D_m(U)-1\bigr)_+,
\]
and
\begin{align}
 \Ee[B(U)-\log L(U)]
 \le{}&
 \sum_{m\le y}\Lambda(m)\Ee D_m(U)
 \notag\\
 &+\frac12\sum_{y<m\le U}\Lambda(m)\Ee(D_m(U))_2.
 \label{eq:abstract-transfer}
\end{align}
\end{theorem}

\begin{proof}
The identities
\[
 \log r=\sum_{m\mid r}\Lambda(m)
\]
and
\[
 \log L(U)=\sum_{m\le U}\Lambda(m)\one_{\{D_m(U)\ge1\}}
\]
give
\[
 B(U)-\log L(U)
 =
 \sum_{m\le U}\Lambda(m)
 \left(D_m(U)-\one_{\{D_m(U)\ge1\}}\right).
\]
The summand in parentheses equals $(D_m(U)-1)_+$.  For every nonnegative
integer $D$,
\[
 (D-1)_+\le D,
 \qquad
 (D-1)_+\le\frac12(D)_2.
\]
Use the first inequality for $m\le y$ and the second for $m>y$.
\end{proof}

\begin{corollary}[Harmonic moment criterion]\label{cor:harmonic-transfer}
Suppose that
\[
 \Ee C_r\le\frac{A}{r}
\]
and
\[
 \Ee\!\left[C_r(C_s-\one_{\{r=s\}})\right]\le\frac{B}{rs}
\]
for all $r,s\le U$.  Then
\[
 \Ee[B(U)-\log L(U)]
 \le
 C\left\{
 A\log(eU)\log(ey)
 +B\frac{\log^2(eU)}{y}
 \right\}.
\]
In particular, with $y=\log^2(eU)$,
\[
 \Ee[B(U)-\log L(U)]
 =
 O\!\left(A\log U\log\log U+B\right).
\]
\end{corollary}

\begin{proof}
The first-moment assumption gives
\[
 \Ee D_m(U)
 \le\frac{A}{m}\sum_{j\le U/m}\frac1j
 \le\frac{A}{m}\log\frac{eU}{m}.
\]
Also,
\[
 (D_m(U))_2
 =
 \sum_{\substack{r,s\le U\\m\mid r,\ m\mid s}}
 C_r(C_s-\one_{\{r=s\}}),
\]
so
\[
 \Ee(D_m(U))_2
 \le
 \frac{B}{m^2}
 \left(\sum_{j\le U/m}\frac1j\right)^2
 \le
 \frac{B}{m^2}\log^2\frac{eU}{m}.
\]
Chebyshev's estimate $\sum_{m\le x}\Lambda(m)=O(x)$ and partial summation
give
\[
 \sum_{m\le y}\frac{\Lambda(m)}m=O(\log(ey)),
 \qquad
 \sum_{m>y}\frac{\Lambda(m)}{m^2}=O(1/y).
\]
Substitution in \eqref{eq:abstract-transfer} proves the result.
\end{proof}

For a process version, put
\[
 B_n(t)=\sum_{r\le n^t}C_{n,r}\log r,
 \qquad
 L_n(t)=\lcm\{r\le n^t:C_{n,r}>0\}.
\]

\begin{lemma}[Functional transfer]\label{lem:functional-transfer}
The function
\[
 t\longmapsto B_n(t)-\log L_n(t)
\]
is nonnegative and nondecreasing.  Consequently,
\[
 \sup_{0\le t\le1}
 \bigl|B_n(t)-\log L_n(t)\bigr|
 = B_n(1)-\log L_n(1).
\]
Thus, if $s_n>0$ and
\[
 \frac{B_n(1)-\log L_n(1)}{s_n}\longrightarrow0
\]
in probability, then the additive and LCM processes differ by
$o_{\Pp}(1)$ uniformly after division by $s_n$.
\end{lemma}

\begin{proof}
When the threshold crosses a length $r$ with $C_{n,r}>0$, the additive
statistic increases by $C_{n,r}\log r$, while the logarithm of the least
common multiple increases by at most $\log r$.  Hence the discrepancy
cannot decrease.
\end{proof}

\section{Exponentially Ewens-stable cycle indices}

We now isolate the probabilistic condition common to standardized
permutations and shelf shuffles.

For a permutation $\pi$, write $C_r(\pi)$ for its number of $r$-cycles.

\begin{definition}[Triangular cycle-index family]\label{def:triangular}
For each target size $n$, let
\[
 \{\Pp_{n,m}:m\ge0\}
\]
be probability measures on $S_m$, with the convention that $S_0$ contains
the empty permutation.  Their cycle index is
\[
 \cZ_n(u,\mathbf x)
 =
 1+\sum_{m\ge1}u^m
 \Ee_{n,m}\prod_{r\ge1}x_r^{C_r}.
\]
We study the diagonal random permutation
\[
 \pi_n\sim\Pp_{n,n}.
\]
\end{definition}

\begin{definition}[Exponential Ewens stability]\label{def:stability}
The triangular family is \emph{exponentially Ewens-stable} if there are
real coefficients $b_{n;r,k}$ and constants $A<\infty$, $\rho\in(0,1)$,
independent of $n$, such that
\begin{equation}\label{eq:cycle-index-form}
 \cZ_n(u,\mathbf x)
 =
 \exp\left\{
 \sum_{r\ge1}\sum_{k\ge1}
 \frac{b_{n;r,k}}{rk}x_r^k u^{rk}
 \right\}
\end{equation}
as a formal power series, and
\begin{align}
 |b_{n;r,1}-1|&\le A\rho^r,
 \label{eq:stability-one}\\
 |b_{n;r,k}|&\le A\rho^{rk},
 \qquad k\ge2.
 \label{eq:stability-high}
\end{align}
The normalization $\cZ_n(u,\mathbf1)=(1-u)^{-1}$ is included in the
definition, since every coefficient at $\mathbf x=\mathbf1$ is one.
\end{definition}

The coefficient $b_{n;r,1}=1$ is the uniform-permutation contribution.
The remaining terms are allowed to depend arbitrarily on $n$, but must be
exponentially confined to short cycle lengths.

Put
\[
 \ell_n=\log n,
 \qquad
 u_{n,r}=\frac{\log r}{\ell_n}.
\]
For a bounded function $f:[0,1]\to\R$, define the logarithmic cycle field
\[
 \cG_n(f)
 =
 \frac1{\sqrt{\ell_n}}
 \sum_{r=1}^n
 f(u_{n,r})\left(C_{n,r}-\frac1r\right).
\]

\begin{theorem}[Gaussian logarithmic cycle field]\label{thm:cycle-field}
Let $(\pi_n)$ be an exponentially Ewens-stable triangular family.  For
every fixed collection of bounded Riemann-integrable functions
$f_1,\ldots,f_d$ on $[0,1]$,
\[
 \bigl(\cG_n(f_1),\ldots,\cG_n(f_d)\bigr)
 \Rightarrow
 \left(
 \int_0^1f_1(u)\,\dd W(u),\ldots,
 \int_0^1f_d(u)\,\dd W(u)
 \right).
\]
Equivalently, the limit is centered Gaussian with covariance
\[
 \Cov(\cG(f),\cG(g))=\int_0^1f(u)g(u)\,\dd u.
\]
\end{theorem}

The main process consequence is the following.

\begin{theorem}[Functional Erd\H{o}s--Tur\'an universality]
\label{thm:general-functional}
Let $(\pi_n)$ be exponentially Ewens-stable, and let $C_{n,r}=C_r(\pi_n)$.
Set
\[
 K_n(t)=\sum_{r\le n^t}C_{n,r},
\]
and
\[
 \cO_n(t)=\log\lcm\{r\le n^t:C_{n,r}>0\}.
\]
Then, in $D([0,1],\R^2)$,
\begin{equation}\label{eq:functional-main}
 \left(
 \frac{K_n(t)-t\ell_n}{\sqrt{\ell_n}},
 \frac{\cO_n(t)-\frac12t^2\ell_n^2}
      {\sqrt{\ell_n^3/3}}
 \right)_{0\le t\le1}
 \Rightarrow
 \left(
 W(t),\sqrt3\int_0^t u\,\dd W(u)
 \right)_{0\le t\le1}.
\end{equation}
The limiting process is continuous and has covariance kernels
\begin{align*}
 \Cov(W(s),W(t))&=\min(s,t),\\
 \Cov\left(
 \sqrt3\int_0^s u\,\dd W(u),
 \sqrt3\int_0^t u\,\dd W(u)
 \right)&=\min(s,t)^3,\\
 \Cov\left(
 W(s),\sqrt3\int_0^t u\,\dd W(u)
 \right)&=\frac{\sqrt3}{2}\min(s,t)^2.
\end{align*}
\end{theorem}

\begin{corollary}[Endpoint joint law]\label{cor:endpoint-joint}
Under the hypotheses of Theorem~\ref{thm:general-functional},
\[
 \left(
 \frac{K_n-\log n}{\sqrt{\log n}},
 \frac{\log\ord(\pi_n)-\frac12(\log n)^2}
      {\sqrt{\frac13(\log n)^3}}
 \right)
 \Rightarrow
 N\!\left(
 0,
 \begin{pmatrix}
 1&\sqrt3/2\\
 \sqrt3/2&1
 \end{pmatrix}
 \right).
\]
In particular,
\[
 \frac{\log\ord(\pi_n)-\frac12(\log n)^2}
 {\sqrt{\frac13(\log n)^3}}
 \Rightarrow N(0,1).
\]
\end{corollary}

\begin{theorem}[Sharp Berry--Esseen universality]
\label{thm:sharp-be}
Let $(\pi_n)$ be an exponentially Ewens-stable triangular family with
stability constants $A$ and $\rho$.  Put
\[
 \ell_n=\log n,
 \qquad
 a_n=\frac12\ell_n^2-\ell_n\log\ell_n,
 \qquad
 s_n=\sqrt{\ell_n^3/3}.
\]
Then there is $C=C(A,\rho)<\infty$ such that, for all $n\ge3$,
\begin{equation}\label{eq:sharp-be}
 \Kol\!\left(
 \frac{\log\ord(\pi_n)-a_n}{s_n},N(0,1)
 \right)
 \le \frac{C}{\sqrt{\ell_n}}.
\end{equation}
The estimate is uniform over every triangular subfamily with common
stability constants.
\end{theorem}

\begin{theorem}[Universal first Edgeworth expansion and optimality]
\label{thm:edgeworth}
Let $(\pi_n)$ be an exponentially Ewens-stable triangular family with
common stability constants $A$ and $\rho$.  Put
\[
 Y_n=\log\ord(\pi_n),\qquad
 m_n=\Ee Y_n,\qquad
 \ell_n=\log n,\qquad
 s_n=\sqrt{\ell_n^3/3}.
\]
Then there is $C=C(A,\rho)<\infty$ such that, for all sufficiently large
$n$,
\begin{equation}\label{eq:edgeworth}
 \sup_{x\in\R}\left|
 \Pp\left\{\frac{Y_n-m_n}{s_n}\le x\right\}
 -\Phi(x)
 -\frac{\sqrt3}{8\sqrt{\ell_n}}(1-x^2)\varphi(x)
 \right|
 \le
 C\left(\frac{\log\ell_n}{\ell_n}\right)^{2/3}.
\end{equation}
Consequently,
\begin{equation}\label{eq:kol-exact}
 \Kol\!\left(\frac{Y_n-m_n}{s_n},N(0,1)\right)
 =
 \frac{\sqrt3}{8\sqrt{2\pi\ell_n}}
 +O_{A,\rho}\!\left(
 \left(\frac{\log\ell_n}{\ell_n}\right)^{2/3}
 \right).
\end{equation}
In particular, the $(\log n)^{-1/2}$ convergence rate is optimal.  Both
statements are uniform over every triangular subfamily with the same
stability constants.
\end{theorem}

The qualitative and quantitative proofs occupy the next sections.

\section{Cumulants from the cycle index}

Fix cutoffs
\[
 b_n=\lceil(\log n)^3\rceil,
 \qquad
 U_n=\left\lfloor\frac{n}{(\log n)^3}\right\rfloor.
\]
For bounded arrays $f_{n,r}^{(j)}$, supported on $b_n<r\le U_n$, put
\[
 S_{n,j}=\sum_{b_n<r\le U_n}f_{n,r}^{(j)}C_{n,r}.
\]

\begin{proposition}[Mixed-cumulant extraction]\label{prop:mixed-cumulant}
Fix $q\ge1$ and $M<\infty$.  There are constants
$C_{q,M}<\infty$ and $\rho_{q,M}\in(0,1)$, depending only on
$q,M,A,\rho$, such that, uniformly over arrays satisfying
\[
 |f_{n,r}^{(j)}|\le M,
 \qquad
 1\le j\le q,
\]
one has
\begin{equation}\label{eq:mixed-cumulant}
 \left|
 \kappa(S_{n,1},\ldots,S_{n,q})
 -
 \sum_{b_n<r\le U_n}
 \frac{\prod_{j=1}^q f_{n,r}^{(j)}}{r}
 \right|
 \le
 C_{q,M}(1+\log n)^q\rho_{q,M}^{\,b_n}.
\end{equation}
\end{proposition}

\begin{proof}
For $\mathbf z=(z_1,\ldots,z_q)$, substitute
\[
 x_r=\exp\left(\sum_{j=1}^qz_jf_{n,r}^{(j)}\right)
\]
for $b_n<r\le U_n$, and $x_r=1$ otherwise.  Since
$\cZ_n(u,\mathbf1)=(1-u)^{-1}$,
\[
 \Ee\exp\left(\sum_{j=1}^qz_jS_{n,j}\right)
 =
 [u^n]\frac1{1-u}\exp H_n(u,\mathbf z),
\]
where
\[
 H_n(u,\mathbf z)
 =
 \sum_{b_n<r\le U_n}\sum_{k\ge1}
 \frac{b_{n;r,k}}{rk}
 \left[
 \exp\left(k\sum_{j=1}^qz_jf_{n,r}^{(j)}\right)-1
 \right]u^{rk}.
\]

For a nonempty set $B\subseteq[q]$, define
\[
 H_{n,B}(u)
 =
 \left.
 \prod_{j\in B}\frac{\partial}{\partial z_j}
 H_n(u,\mathbf z)
 \right|_{\mathbf z=0}.
\]
Then
\[
 H_{n,B}(u)
 =
 \sum_{b_n<r\le U_n}
 \frac{\prod_{j\in B}f_{n,r}^{(j)}}r
 \left[
 b_{n;r,1}u^r+
 \sum_{k\ge2}b_{n;r,k}k^{|B|-1}u^{rk}
 \right].
\]
Write
\[
 H_{n,B}(u)=A_{n,B}(u)+E_{n,B}(u),
\]
where
\[
 A_{n,B}(u)
 =
 \sum_{b_n<r\le U_n}
 \frac{\prod_{j\in B}f_{n,r}^{(j)}}r u^r.
\]
By \eqref{eq:stability-one}--\eqref{eq:stability-high}, the coefficients
of $E_{n,B}$ have exponentially summable absolute values, and
\begin{equation}\label{eq:E-small-l1}
 \sum_{m\ge0}|[u^m]E_{n,B}(u)|
 \le C_{q,M}\rho_1^{b_n}
\end{equation}
for some $\rho_1\in(0,1)$.

The derivative
\[
 \left.
 \partial_{z_1}\cdots\partial_{z_q}
 e^{H_n(u,\mathbf z)}
 \right|_{\mathbf z=0}
\]
is the sum, over set partitions $\pi$ of $[q]$, of
\[
 \prod_{B\in\pi}H_{n,B}(u).
\]
The product of the $A_{n,B}$ terms has degree at most $qU_n<n$ for all
large $n$.  Therefore
\[
 [u^n]\frac1{1-u}\prod_{B\in\pi}A_{n,B}(u)
 =
 \prod_{B\in\pi}A_{n,B}(1).
\]
Every remaining product contains at least one $E_{n,B}$, and its total
absolute coefficient mass is
\[
 O_{q,M}\bigl((1+\log n)^q\rho_1^{b_n}\bigr)=o(1)
\]
by \eqref{eq:E-small-l1}.  It follows that all joint moments of order at
most $q$ differ by $o(1)$ from those of the formal random vector with
cumulant generating function $H_n(1,\mathbf z)$.

Cumulants are universal polynomials in moments.  Hence
\[
 \kappa(S_{n,1},\ldots,S_{n,q})
 =
 \left.
 \partial_{z_1}\cdots\partial_{z_q}
 H_n(1,\mathbf z)
 \right|_{\mathbf z=0}
 +o(1).
\]
The derivative on the right is
\[
 \sum_{b_n<r\le U_n}
 \frac{\prod_jf_{n,r}^{(j)}}r
 \left[
 b_{n;r,1}
 +\sum_{k\ge2}b_{n;r,k}k^{q-1}
 \right].
\]
The difference from the main term in \eqref{eq:mixed-cumulant} is bounded
by $C_{q,M}\rho_1^{b_n}$, again using exponential stability.  Combining
the two bounds and enlarging the constants gives the displayed estimate.
\end{proof}

The same cycle index gives the moment estimates needed for the arithmetic
transfer.

\begin{proposition}[Harmonic one- and two-cycle bounds]
\label{prop:harmonic-moments}
There is a constant $C<\infty$, depending only on $A,\rho$, such that,
for all $n$ and all $r,s\le n$,
\[
 \Ee C_{n,r}\le\frac{C}{r},
\]
and
\[
 \Ee\!\left[
 C_{n,r}(C_{n,s}-\one_{\{r=s\}})
 \right]
 \le\frac{C}{rs}.
\]
\end{proposition}

\begin{proof}
Differentiate \eqref{eq:cycle-index-form} with respect to $x_r$ and set
$\mathbf x=\mathbf1$.  Since $\cZ_n(u,\mathbf1)=(1-u)^{-1}$,
\[
 \Ee C_{n,r}
 =
 [u^n]\frac1{1-u}
 \left(
 \frac1r\sum_{k\ge1}b_{n;r,k}u^{rk}
 \right).
\]
Thus
\[
 |\Ee C_{n,r}|
 \le\frac1r\sum_{k\ge1}|b_{n;r,k}|
 \le\frac{C}{r}.
\]

For $r\ne s$, two differentiations give
\[
 \Ee(C_{n,r}C_{n,s})
 =
 [u^n]\frac1{1-u}A_{n,r}(u)A_{n,s}(u),
\]
where
\[
 A_{n,r}(u)=\frac1r\sum_{k\ge1}b_{n;r,k}u^{rk}.
\]
The absolute coefficient mass of $A_{n,r}$ is at most $C/r$, proving the
claim.

For $r=s$,
\[
 \Ee(C_{n,r})_2
 =
 [u^n]\frac1{1-u}
 \left(A_{n,r}(u)^2+B_{n,r}(u)\right),
\]
where
\[
 B_{n,r}(u)
 =
 \frac1r\sum_{k\ge2}(k-1)b_{n;r,k}u^{rk}.
\]
The first term has absolute coefficient mass $O(r^{-2})$.  The second has
mass
\[
 O\left(\frac{\rho^{2r}}r\right)=O(r^{-2}),
\]
because $\sup_{r\ge1}r\rho^{2r}<\infty$.
\end{proof}

\section{Long-cycle coupling, sharp Berry--Esseen, and Edgeworth transfer}
\label{sec:be}

The exponential coefficient condition gives substantially more than fixed
cumulant control: after a growing short-cycle cutoff, the complete remaining
cycle vector is close in total variation to that of a uniform permutation.
This permits direct transfer of the sharp classical rate and of the first classical Edgeworth term.

Let $\Pi_n$ be uniform on $S_n$, and write $C^0_{n,r}=C_r(\Pi_n)$.
For an integer $b\ge1$, set
\[
 \mathbf C_n^{>b}=(C_{n,r})_{r>b},
 \qquad
 \mathbf C_n^{0,>b}=(C^0_{n,r})_{r>b}.
\]

\begin{proposition}[Exponential long-cycle total-variation universality]
\label{prop:long-cycle-tv}
For every exponentially Ewens-stable triangular family, there is
$C=C(A,\rho)<\infty$ such that, uniformly in $n$ and $b\ge1$,
\begin{equation}\label{eq:long-cycle-tv}
 d_{\TV}\!\left(
 \cL(\mathbf C_n^{>b}),
 \cL(\mathbf C_n^{0,>b})
 \right)
 \le C\rho^b.
\end{equation}
\end{proposition}

\begin{proof}
Let
\[
 d_{n;r,k}=b_{n;r,k}-\one_{\{k=1\}}.
\]
The uniform-permutation cycle index is
\[
 \cZ_0(u,\mathbf x)
 =\exp\left\{\sum_{r\ge1}\frac{x_ru^r}{r}\right\}.
\]
Since both cycle indices equal $(1-u)^{-1}$ at $\mathbf x=\mathbf1$,
\[
 \exp\left\{
 \sum_{r,k\ge1}\frac{d_{n;r,k}}{rk}u^{rk}
 \right\}=1.
\]
The exponent has zero constant term, so it vanishes as a formal power
series.  After setting $x_r=1$ for $r\le b$, we therefore have the exact
factorization
\begin{equation}\label{eq:long-index-factor}
 \cZ_n^{>b}(u,\mathbf x)
 =\cZ_0^{>b}(u,\mathbf x)\exp R_{n,b}(u,\mathbf x),
\end{equation}
where
\[
 R_{n,b}(u,\mathbf x) =\sum_{r>b}\sum_{k\ge1}
 \frac{d_{n;r,k}}{rk}(x_r^k-1)u^{rk}.
\]

For an absolutely summable multivariate series $F$, let $\|F\|_1$ denote
its total absolute coefficient mass.  Exponential stability gives
\begin{align*}
 \|R_{n,b}\|_1
 &\le
 2A\sum_{r>b}\frac{\rho^r}{r}
 +2A\sum_{r>b}\sum_{k\ge2}\frac{\rho^{rk}}{rk}
 \\
 &\le C\rho^b.
\end{align*}
The coefficient $\ell^1$ norm is submultiplicative, and hence
\begin{equation}\label{eq:exp-R-l1}
 \|e^{R_{n,b}}-1\|_1
 \le e^{\|R_{n,b}\|_1}-1
 \le C\rho^b.
\end{equation}

For each $m\ge0$, the coefficient polynomial
$[u^m]\cZ_0^{>b}$ has nonnegative coefficients whose sum is
\[
 [u^m]\cZ_0^{>b}(u,\mathbf1)
 =[u^m](1-u)^{-1}=1.
\]
Taking the coefficient of $u^n$ in
\eqref{eq:long-index-factor}, using convolution and
\eqref{eq:exp-R-l1}, gives
\[
 \sum_{\mathbf c}
 \left|
 \Pp\{\mathbf C_n^{>b}=\mathbf c\}
 -\Pp\{\mathbf C_n^{0,>b}=\mathbf c\}
 \right|
 \le C\rho^b.
\]
Division by two proves \eqref{eq:long-cycle-tv}.
\end{proof}

We import the sharp classical endpoint theorem in the following form.

\begin{inputtheorem}[Barbour--Tavar\'e sharp rate
\cite{BarbourTavare1994,Zacharovas2009}]
\label{input:BT}
Let
\[
 Y_n^0=\log\ord(\Pi_n),
 \qquad
 \ell_n=\log n,
 \qquad
 a_n=\frac12\ell_n^2-\ell_n\log\ell_n,
 \qquad
 s_n=\sqrt{\ell_n^3/3}.
\]
Then
\begin{equation}\label{eq:BT}
 \Kol\!\left(\frac{Y_n^0-a_n}{s_n},N(0,1)\right)
 \le \frac{C_0}{\sqrt{\ell_n}}.
\end{equation}
\end{inputtheorem}

\begin{inputtheorem}[Zacharovas first Edgeworth expansion
\cite{Zacharovas2009}]
\label{input:Zacharovas}
Let
\[
 Y_n^0=\log\ord(\Pi_n),\qquad
 m_n^0=\Ee Y_n^0,\qquad
 \ell_n=\log n,\qquad
 s_n=\sqrt{\ell_n^3/3}.
\]
Then, for all sufficiently large $n$,
\begin{equation}\label{eq:Zacharovas}
 \sup_{x\in\R}\left|
 \Pp\left\{\frac{Y_n^0-m_n^0}{s_n}\le x\right\}
 -\Phi(x)
 -\frac{\sqrt3}{8\sqrt{\ell_n}}(1-x^2)\varphi(x)
 \right|
 \le
 C_1\left(\frac{\log\ell_n}{\ell_n}\right)^{2/3}.
\end{equation}
\end{inputtheorem}

The following elementary lemmas record the quantitative coupling transfer.

\begin{lemma}[Kolmogorov transfer under a partial coupling]
\label{lem:kol-coupling}
Let $X,X_0$ be coupled real random variables.  Suppose there are an event
$G$, a nonnegative random variable $\Delta$, and numbers $\eta,d\ge0$ such
that
\[
 \Pp(G^c)\le\eta,
 \qquad
 |X-X_0|\le\Delta\quad\hbox{on }G,
 \qquad
 \Ee\Delta\le d.
\]
Then, for every $a\in\R$, $s>0$, and $h>0$,
\begin{equation}\label{eq:kol-coupling}
 \Kol\!\left(\frac{X-a}{s},N(0,1)\right)
 \le
 \Kol\!\left(\frac{X_0-a}{s},N(0,1)\right)
 +\eta+C\frac{h}{s}+\frac{d}{h}.
\end{equation}
\end{lemma}

\begin{proof}
For every $z\in\R$,
\[
 \Pp\{X\le z\}
 \le
 \Pp\{X_0\le z+h\}+\Pp(G^c)+\Pp\{\Delta>h\},
\]
and the analogous lower bound holds with $z-h$.  The standard normal
cdf has density bounded by $(2\pi)^{-1/2}$, while Markov's inequality gives
$\Pp\{\Delta>h\}\le d/h$.  Taking the supremum over
$z=a+sx$ proves the claim.
\end{proof}

\begin{lemma}[Distributional transfer under a partial coupling]
\label{lem:cdf-coupling}
Under the hypotheses of Lemma~\ref{lem:kol-coupling}, for every $h>0$,
\begin{equation}\label{eq:cdf-coupling}
 \sup_{z\in\R}|\Pp\{X\le z\}-\Pp\{X_0\le z\}|
 \le
 \eta+\frac{d}{h}
 +\sup_{z\in\R}\Pp\{z-h<X_0\le z+h\}.
\end{equation}
\end{lemma}

\begin{proof}
The upper bound
\[
 \Pp\{X\le z\}
 \le
 \Pp\{X_0\le z+h\}+\eta+\Pp\{\Delta>h\}
\]
and the corresponding lower bound with $z-h$ imply
\eqref{eq:cdf-coupling} after Markov's inequality.
\end{proof}

\begin{proof}[Proof of Theorem~\ref{thm:sharp-be}]
Put
\[
 Y_n=\log\ord(\pi_n).
\]
Choose
\[
 b_n=
 \left\lceil\frac{2\log\ell_n}{|\log\rho|}\right\rceil
\]
when $\ell_n>1$, and absorb the finitely many remaining $n$ into the
constant.  By Proposition~\ref{prop:long-cycle-tv}, the long-cycle vectors
can be coupled so that
\begin{equation}\label{eq:coupling-failure}
 \Pp\{\mathbf C_n^{>b_n}\ne\mathbf C_n^{0,>b_n}\}
 \le C\rho^{b_n}\le C\ell_n^{-2}.
\end{equation}

Let $Y_n^{>b_n}$ and $Y_n^{0,>b_n}$ be the logarithms of the least common
multiples of the cycle lengths exceeding $b_n$ in the two permutations.
On the successful coupling event these variables are equal.  Moreover,
\begin{align*}
 0\le Y_n-Y_n^{>b_n}
 &\le S_{n,b_n}:=\sum_{r\le b_n}C_{n,r}\log r,\\
 0\le Y_n^0-Y_n^{0,>b_n}
 &\le S^0_{n,b_n}:=\sum_{r\le b_n}C^0_{n,r}\log r.
\end{align*}
Hence, on that event,
\[
 |Y_n-Y_n^0|\le S_{n,b_n}+S^0_{n,b_n}.
\]
Proposition~\ref{prop:harmonic-moments}, together with
$\Ee C^0_{n,r}=1/r$, yields
\begin{equation}\label{eq:short-cycle-mean}
 \Ee(S_{n,b_n}+S^0_{n,b_n})
 \le C\sum_{r\le b_n}\frac{\log r}{r}
 \le C\log^2(e b_n).
\end{equation}

Apply Lemma~\ref{lem:kol-coupling} with $X=Y_n$, $X_0=Y_n^0$,
$h=\ell_n$, and the center and scale from
Input Theorem~\ref{input:BT}.  Equations
\eqref{eq:BT}--\eqref{eq:short-cycle-mean} give
\begin{align*}
 \Kol\!\left(\frac{Y_n-a_n}{s_n},N(0,1)\right)
 &\le
 \frac{C}{\sqrt{\ell_n}}
 +C\ell_n^{-2}
 +C\frac{\ell_n}{\ell_n^{3/2}}
 +C\frac{\log^2(e b_n)}{\ell_n}\\
 &\le \frac{C}{\sqrt{\ell_n}},
\end{align*}
which proves \eqref{eq:sharp-be}.
\end{proof}

\begin{proof}[Proof of Theorem~\ref{thm:edgeworth}]
Let $Y_n^0=\log\ord(\Pi_n)$ and $m_n^0=\Ee Y_n^0$.  Choose a fixed
$c>0$ so large that
\[
 c|\log\rho|>4,
\]
and put
\[
 b_n=\lceil c\ell_n\rceil.
\]
By Proposition~\ref{prop:long-cycle-tv}, the long-cycle vectors can be
maximally coupled so that, for an event $G_n$,
\begin{equation}\label{eq:edge-coupling-failure}
 \Pp(G_n^c)\le C\rho^{b_n}=O(n^{-4}).
\end{equation}
After this maximal coupling, sample the short-cycle coordinates according
to their respective conditional laws given the long vectors.  This
extends the coupling to the complete cycle vectors without changing either
marginal law.

On $G_n$, the least common multiples of the cycle lengths exceeding $b_n$
are identical.  Therefore
\begin{equation}\label{eq:edge-coupling-size}
 |Y_n-Y_n^0|
 \le
 \Delta_n:=S_{n,b_n}+S_{n,b_n}^0,
\end{equation}
where
\[
 S_{n,b}=
 \sum_{r\le b}C_{n,r}\log r,
 \qquad
 S_{n,b}^0=
 \sum_{r\le b}C_{n,r}^0\log r.
\]
Proposition~\ref{prop:harmonic-moments} and
$\Ee C_{n,r}^0=1/r$ give
\begin{equation}\label{eq:edge-short-mean}
 d_n:=\Ee\Delta_n
 \le C\sum_{r\le b_n}\frac{\log r}{r}
 \le C\log^2(e b_n)
 =O((\log\ell_n)^2).
\end{equation}

We first compare the two laws with the common center $m_n^0$.  Let
\[
 \widetilde F_n(x)=
 \Pp\left\{\frac{Y_n-m_n^0}{s_n}\le x\right\},
 \qquad
 F_n^0(x)=
 \Pp\left\{\frac{Y_n^0-m_n^0}{s_n}\le x\right\}.
\]
Set
\[
 h_n=\ell_n^{3/4},
 \qquad
 r_n=\left(\frac{\log\ell_n}{\ell_n}\right)^{2/3}.
\]
Lemma~\ref{lem:cdf-coupling} gives
\begin{align}
 \sup_x|\widetilde F_n(x)-F_n^0(x)|
 \le{}& O(n^{-4})+\frac{d_n}{h_n}
 \notag\\
 &+\sup_z\Pp\{z-h_n<Y_n^0\le z+h_n\}.
 \label{eq:edge-cdf-transfer}
\end{align}

Define the Edgeworth approximant
\[
 E_n(x)=
 \Phi(x)+
 \frac{\sqrt3}{8\sqrt{\ell_n}}(1-x^2)\varphi(x).
\]
The functions $E_n$ are uniformly Lipschitz: indeed,
\[
 E_n'(x)=
 \varphi(x)+
 \frac{\sqrt3}{8\sqrt{\ell_n}}(x^3-3x)\varphi(x),
\]
and the polynomial times Gaussian factor is bounded.  Input
Theorem~\ref{input:Zacharovas} therefore yields
\begin{equation}\label{eq:uniform-anti-edge}
 \sup_z\Pp\{z-h_n<Y_n^0\le z+h_n\}
 \le
 C\frac{h_n}{s_n}+Cr_n.
\end{equation}
Since
\[
 \frac{d_n}{h_n}
 =O\!\left(\frac{(\log\ell_n)^2}{\ell_n^{3/4}}\right)
 =o(r_n),
 \qquad
 \frac{h_n}{s_n}=O(\ell_n^{-3/4})=o(r_n),
\]
equations \eqref{eq:edge-cdf-transfer} and
\eqref{eq:uniform-anti-edge} give
\begin{equation}\label{eq:edge-common-center}
 \sup_x|\widetilde F_n(x)-F_n^0(x)|=O_{A,\rho}(r_n).
\end{equation}

It remains to replace the uniform mean by the actual mean.  For every
permutation of $n$ letters,
\[
 0\le\log\ord(\pi)
 \le\sum_{C\in\Cyc(\pi)}\log|C|
 \le\sum_{C\in\Cyc(\pi)}|C|=n.
\]
Hence, using \eqref{eq:edge-coupling-failure},
\eqref{eq:edge-coupling-size}, and \eqref{eq:edge-short-mean},
\begin{equation}\label{eq:mean-transfer-edge}
 |m_n-m_n^0|
 \le d_n+2n\Pp(G_n^c)
 =O((\log\ell_n)^2).
\end{equation}
Put
\[
 \gamma_n=\frac{m_n-m_n^0}{s_n}.
\]
Then $|\gamma_n|=O((\log\ell_n)^2/\ell_n^{3/2})=o(r_n)$.  Since
\[
 \Pp\left\{\frac{Y_n-m_n}{s_n}\le x\right\}
 =\widetilde F_n(x+\gamma_n),
\]
Input Theorem~\ref{input:Zacharovas},
\eqref{eq:edge-common-center}, and the uniform Lipschitz bound for $E_n$
prove \eqref{eq:edgeworth}.

Finally, elementary calculus gives
\[
 \sup_{x\in\R}|(1-x^2)\varphi(x)|=\varphi(0)=\frac1{\sqrt{2\pi}}.
\]
Taking suprema in \eqref{eq:edgeworth} proves
\eqref{eq:kol-exact} and hence optimality.
\end{proof}

\begin{remark}[Why the adjusted center is necessary]
The qualitative theorem may be centered at $\frac12\ell_n^2$ because
$\ell_n\log\ell_n=o(\ell_n^{3/2})$.  At Berry--Esseen resolution this is
no longer harmless: changing from $a_n$ to $\frac12\ell_n^2$ translates
the normalized variable by
\[
 \frac{\ell_n\log\ell_n}{s_n}
 =\sqrt3\,\frac{\log\ell_n}{\sqrt{\ell_n}},
\]
which is larger by a factor $\log\ell_n$ than the sharp
$\ell_n^{-1/2}$ error.  Thus the Barbour--Tavar\'e correction is intrinsic
to the sharp formulation.
\end{remark}

\section{The Gaussian cycle field and functional tightness}

We first record the logarithmic Riemann sums used throughout.

\begin{lemma}[Logarithmic Riemann sums]\label{lem:riemann}
If $f$ is bounded and Riemann integrable on $[0,1]$, then
\[
 \frac1{\log n}
 \sum_{r=1}^n\frac{f(\log r/\log n)}r
 \longrightarrow
 \int_0^1f(u)\,\dd u.
\]
For every fixed continuous $f$, the corresponding truncated sums up to
$n^t$ converge uniformly in $t\in[0,1]$:
\[
 \sup_{0\le t\le1}
 \left|
 \frac1{\log n}
 \sum_{r\le n^t}\frac{f(\log r/\log n)}r
 -
 \int_0^t f(u)\,\dd u
 \right|
 \longrightarrow0.
\]
\end{lemma}

\begin{proof}
Put $u_{n,r}=\log r/\log n$.  The mesh
$\max_r(u_{n,r+1}-u_{n,r})$ is $O((\log n)^{-1})$, and
\[
 u_{n,r+1}-u_{n,r}
 =
 \frac1{\log n}\log\left(1+\frac1r\right)
 =
 \frac1{r\log n}+O\left(\frac1{r^2\log n}\right).
\]
Thus the displayed sum differs by $o(1)$ from the Riemann sum associated
with the partition $(u_{n,r})$.  This proves the first assertion for
continuous functions.  Bounded Riemann-integrable functions follow by
upper and lower step-function approximation.  For continuous $f$, the
same comparison applied to every initial segment of the partition is
uniform in its endpoint, proving the second assertion.
\end{proof}

\begin{proof}[Proof of Theorem~\ref{thm:cycle-field}]
First replace the full sums by
\[
 \cG_n^\circ(f)
 =
 \frac1{\sqrt{\ell_n}}
 \sum_{b_n<r\le U_n}
 f(u_{n,r})\left(C_{n,r}-\frac1r\right).
\]
Proposition~\ref{prop:harmonic-moments} gives
\[
 \Ee\sum_{r\le b_n}C_{n,r}=O(\log b_n)=O(\log\log n)
\]
and
\[
 \Ee\sum_{r>U_n}C_{n,r}
 =O(\log(n/U_n)+1)=O(\log\log n).
\]
The corresponding deterministic harmonic sums have the same orders.
Hence
\[
 \cG_n(f)-\cG_n^\circ(f)\longrightarrow0
\]
in $L^1$, uniformly for $f$ in a fixed bounded set.

For fixed $f_1,\ldots,f_d$, apply
Proposition~\ref{prop:mixed-cumulant}.  The second joint cumulants satisfy
\[
 \kappa\bigl(\cG_n^\circ(f_i),\cG_n^\circ(f_j)\bigr)
 =
 \frac1{\ell_n}
 \sum_{b_n<r\le U_n}
 \frac{f_i(u_{n,r})f_j(u_{n,r})}{r}
 +o(1),
\]
which converges to $\int_0^1f_i f_j$ by
Lemma~\ref{lem:riemann}.  A joint cumulant of order $q\ge3$ is
\[
 O(\ell_n^{1-q/2})+o(1),
\]
and therefore tends to zero.  The first cumulants tend to zero, the second cumulants have the displayed
limits, and every fixed higher cumulant tends to zero.  Hence all moments
of every fixed linear combination converge to the corresponding Gaussian
moments.  Since the Gaussian law is moment-determinate, the method of
moments and Cram\'er--Wold give the asserted finite-dimensional limit.
\end{proof}

We next prove the process theorem first for the additive cycle field.
Define
\[
 B_n(t)=\sum_{r\le n^t}(\log r)C_{n,r}.
\]

\begin{proposition}[Bivariate additive functional limit]
\label{prop:additive-functional}
In $D([0,1],\R^2)$,
\[
 \left(
 \frac{K_n(t)-t\ell_n}{\sqrt{\ell_n}},
 \frac{B_n(t)-\frac12t^2\ell_n^2}{\sqrt{\ell_n^3/3}}
 \right)_{0\le t\le1}
 \Rightarrow
 \left(
 W(t),\sqrt3\int_0^t u\,\dd W(u)
 \right)_{0\le t\le1}.
\]
\end{proposition}

\begin{proof}
We first work on the intermediate window and center by exact expectations.
Put
\[
 \mathbf X_n^\circ(t)
 =
 \frac1{\sqrt{\ell_n}}
 \sum_{\substack{b_n<r\le U_n\\r\le n^t}}
 \begin{pmatrix}
 1\\
 \sqrt3\,u_{n,r}
 \end{pmatrix}
 \bigl(C_{n,r}-\Ee C_{n,r}\bigr).
\]

Finite-dimensional convergence follows from
Proposition~\ref{prop:mixed-cumulant} and
Lemma~\ref{lem:riemann}.  Indeed, the covariance matrix at times $s,t$
converges to
\[
 \int_0^{\min(s,t)}
 \begin{pmatrix}
 1\\ \sqrt3u
 \end{pmatrix}
 \begin{pmatrix}
 1&\sqrt3u
 \end{pmatrix}\dd u.
\]
All joint cumulants of order at least three vanish after normalization.

We verify tightness.  Consider two adjacent increments over $(s,t]$ and
$(t,u]$, and any fixed scalar projections $A_n,B_n$ of these increments
before division by $\sqrt{\ell_n}$.  Their length supports are disjoint.
The moment--cumulant identity gives
\[
 \Ee(A_n^2B_n^2)
 =
 \kappa(A_n,A_n,B_n,B_n)
 +\kappa(A_n,A_n)\kappa(B_n,B_n)
 +2\kappa(A_n,B_n)^2.
\]By Proposition~\ref{prop:mixed-cumulant}, the main terms in the two mixed
cumulants vanish because the supports are disjoint.  Moreover, for a
deterministic sequence $\eta_n\to0$, independent of $s,t,u$,
\[
 \kappa(A_n,A_n)
 \le C\sum_{\substack{n^s<r\le n^t\\b_n<r\le U_n}}\frac1r+\eta_n,
\]
and similarly for $B_n$, while the absolute values of the mixed second and
fourth cumulants are at most $\eta_n$.  Integral comparison gives
\[
 \sum_{\substack{n^s<r\le n^t\\b_n<r\le U_n}}\frac1r
 \le \ell_n(t-s)+O(b_n^{-1}).
\]
After division by $\ell_n^2$ and use of
$(t-s)(u-t)\le (u-s)^2/4$, we obtain another deterministic
$\widetilde\eta_n\to0$ such that
\[
 \Ee\left[
 \left|\frac{A_n}{\sqrt{\ell_n}}\right|^2
 \left|\frac{B_n}{\sqrt{\ell_n}}\right|^2
 \right]
 \le C(u-s)^2+\widetilde\eta_n
\]
uniformly in $0\le s\le t\le u\le1$.  Lemma~\ref{lem:tightness-criterion},
the standard vanishing-error form of Billingsley's two-increment
criterion, now yields tightness of each coordinate and hence of the
bivariate process.

For completeness, the maximal jump also vanishes.  By
Proposition~\ref{prop:harmonic-moments},
\[
 \Pp\left\{\exists r>b_n:C_{n,r}\ge2\right\}
 \le\frac12\sum_{r>b_n}\Ee(C_{n,r})_2
 =O(b_n^{-1}).
\]
With probability tending to one, every intermediate jump in either
normalized coordinate is therefore $O(\ell_n^{-1/2})$.  Every subsequential
limit is continuous.

It remains to restore the omitted lengths and replace exact means by the
classical centers.  Proposition~\ref{prop:harmonic-moments} gives
\begin{align*}
 \Ee\sum_{r\le b_n}C_{n,r}
 +\Ee\sum_{r>U_n}C_{n,r}
 &=O(\log\log n),\\
 \Ee\sum_{r\le b_n}(\log r)C_{n,r}
 +\Ee\sum_{r>U_n}(\log r)C_{n,r}
 &=O(\ell_n\log\log n).
\end{align*}
Since these omitted processes are nondecreasing, the same bounds control
their suprema in $t$.  They are negligible after the stated
normalizations.

Finally, Proposition~\ref{prop:mixed-cumulant} with $q=1$ and the bounded
weights
$\one_{\{r\le n^t\}}$ and
$u_{n,r}\one_{\{r\le n^t\}}$ gives, uniformly in $t$,
\begin{align*}
 \Ee K_n(t)&=t\ell_n+O(\log\log n),\\
 \Ee B_n(t)&=\frac12t^2\ell_n^2+O(\ell_n\log\log n).
\end{align*}
This completes the proof.
\end{proof}

\section{Proof of the functional order theorem}

\begin{proof}[Proof of Theorem~\ref{thm:general-functional}]
Proposition~\ref{prop:harmonic-moments} and
Corollary~\ref{cor:harmonic-transfer}, with $U=n$, give
\[
 \Ee\left[
 B_n(1)-\log\ord(\pi_n)
 \right]
 =
 O(\ell_n\log\ell_n).
\]
Indeed, choose $y=\ell_n^2$.  The constants in the harmonic moment bounds
are uniform over the triangular family.

By Lemma~\ref{lem:functional-transfer},
\[
 \sup_{0\le t\le1}
 \left|B_n(t)-\cO_n(t)\right|
 =
 B_n(1)-\log\ord(\pi_n).
\]
Consequently,
\[
 \sup_{0\le t\le1}
 \frac{|B_n(t)-\cO_n(t)|}{\ell_n^{3/2}}
 \longrightarrow0
\]
in $L^1$.  Apply Proposition~\ref{prop:additive-functional} and Slutsky's
theorem in $D([0,1],\R^2)$.
\end{proof}

\section{Random standardized permutations}

Let $I_n$ be a finite or countable totally ordered alphabet, and let
\[
 G_{n,1},\ldots,G_{n,m}
\]
be i.i.d.\ with probability vector
\[
 p^{(n)}=(p_{n,a})_{a\in I_n}.
\]
For a word $g=(g_1,\ldots,g_m)$, its standardization is the permutation
$\operatorname{std}(g)\in S_m$ obtained by ranking the letters increasingly,
breaking ties from left to right.  We apply the same letter distribution
$p^{(n)}$ at every auxiliary size $m$, and take the diagonal permutation
\[
 \sigma_n=\operatorname{std}(G_{n,1},\ldots,G_{n,n}).
\]

A word is primitive if it is not a nontrivial power.  Two primitive words
are conjugate if one is a cyclic rotation of the other.  Let
$\widetilde\cQ_r(I_n)$ be the set of conjugacy classes of primitive words
of length $r$.  For $\alpha\in\widetilde\cQ_r(I_n)$, put
\[
 p_\alpha=\prod_{j=1}^r p_{n,\alpha_j};
\]
this is independent of the representative.

We use the following exact result of Guerder.

\begin{inputtheorem}[Primitive-necklace cycle counts
\cite{Guerder2026}]\label{input:guerder}
For each primitive-necklace class $\alpha$, let $D_\alpha$ be the number
of cycles of $\sigma_n$ of type $\alpha$.  If
$\alpha_1,\ldots,\alpha_s$ are distinct classes, of respective lengths
$r_1,\ldots,r_s$, then for all $\ell_1,\ldots,\ell_s\ge0$,
\[
 \Pp(D_{\alpha_1}\ge\ell_1,\ldots,D_{\alpha_s}\ge\ell_s)
 =
 \begin{cases}
 \prod_{j=1}^s p_{\alpha_j}^{\ell_j},
 &\sum_jr_j\ell_j\le n,\\
 0,&\sum_jr_j\ell_j>n.
 \end{cases}
\]
Moreover,
\[
 C_{n,r}=\sum_{\alpha\in\widetilde\cQ_r(I_n)}D_\alpha.
\]
\end{inputtheorem}

\begin{proposition}[Standardized-permutation cycle index]
\label{prop:standardized-index}
For the auxiliary standardized permutations generated by $p^{(n)}$,
\[
 \cZ_n(u,\mathbf x)
 =
 \prod_{r\ge1}
 \prod_{\alpha\in\widetilde\cQ_r(I_n)}
 \frac1{1-p_\alpha x_r u^r}.
\]
Equivalently,
\[
 \cZ_n(u,\mathbf x)
 =
 \exp\left\{
 \sum_{r\ge1}\sum_{k\ge1}
 \frac{b_{n;r,k}}{rk}x_r^k u^{rk}
 \right\},
\]
where
\begin{equation}\label{eq:standardized-b}
 b_{n;r,k}
 =
 r\sum_{\alpha\in\widetilde\cQ_r(I_n)}p_\alpha^k
 =
 \sum_{w\in\cQ_r(I_n)}p_w^k.
\end{equation}
\end{proposition}

\begin{proof}
Let $N$ be independent with
\[
 \Pp(N=m)=(1-u)u^m,\qquad m\ge0.
\]
Mixing the identity in Input Theorem~\ref{input:guerder} over $N$ gives
\[
 \Pp(D_{\alpha_1}\ge\ell_1,\ldots,D_{\alpha_s}\ge\ell_s)
 =
 \prod_{j=1}^s(p_{\alpha_j}u^{r_j})^{\ell_j}.
\]
Thus, under geometric size, the $D_\alpha$ are independent geometric
variables with tails
\[
 \Pp(D_\alpha\ge\ell)=(p_\alpha u^{|\alpha|})^\ell.
\]
Their joint probability generating function is
\[
 (1-u)\cZ_n(u,\mathbf x)
 =
 \prod_\alpha
 \frac{1-p_\alpha u^{|\alpha|}}
      {1-p_\alpha x_{|\alpha|}u^{|\alpha|}}.
\]
Since the total number of letters equals
$\sum_\alpha|\alpha|D_\alpha=N$, setting all $x_r=0$ gives the weighted
Witt identity
\[
 \prod_\alpha(1-p_\alpha u^{|\alpha|})=1-u.
\]
The first displayed product follows.  Taking logarithms and using that a
primitive conjugacy class of length $r$ has exactly $r$ elements gives
\eqref{eq:standardized-b}.
\end{proof}

\begin{lemma}[M\"obius formula and exponential stability]
\label{lem:standardized-stability}
For every $r,k\ge1$,
\begin{equation}\label{eq:mobius-b}
 b_{n;r,k}
 =
 \sum_{d\mid r}\mu(d)
 \left(\sum_{a\in I_n}p_{n,a}^{kd}\right)^{r/d}.
\end{equation}
If
\[
 \sup_n\|p^{(n)}\|_\infty\le R<1,
\]
then there are $A<\infty$ and $\rho\in(0,1)$, depending only on $R$, such
that
\[
 |b_{n;r,1}-1|\le A\rho^r,
 \qquad
 |b_{n;r,k}|\le A\rho^{rk}\quad(k\ge2).
\]
\end{lemma}

\begin{proof}
Every word has a unique primitive root.  M\"obius inversion on the divisor
lattice gives \eqref{eq:mobius-b}.

Put
\[
 s_j=\sum_a p_{n,a}^j.
\]
Since $\max_ap_{n,a}\le R$,
\[
 s_j\le R^{j-1}.
\]
For $k=1$, the $d=1$ term in \eqref{eq:mobius-b} equals one.  Every term
with $d\ge2$ has absolute value at most
\[
 R^{(d-1)r/d}\le R^{r/2}.
\]
For $k\ge2$, every divisor term is bounded by
\[
 R^{(kd-1)r/d}=R^{kr-r/d}\le R^{kr/2}.
\]
There are at most $r$ divisors.  Choose any
\[
 \sqrt R<\rho<1.
\]
The polynomial divisor factor is absorbed into $A\rho^{rk}$.
\end{proof}

Combining the general theory with
Lemma~\ref{lem:standardized-stability} proves the main standardized result.

\begin{theorem}[Standardized-permutation universality]
\label{thm:standardized}
Let
\[
 \sigma_n=\operatorname{std}(G_{n,1},\ldots,G_{n,n}),
\]
where the letters are i.i.d.\ with an arbitrary finite or countable
ordered-alphabet distribution $p^{(n)}$.  Suppose
\[
 \sup_n\|p^{(n)}\|_\infty\le R<1.
\]
Then Theorems~\ref{thm:cycle-field},
\ref{thm:general-functional}, \ref{thm:sharp-be}, and~\ref{thm:edgeworth},
and Corollary~\ref{cor:endpoint-joint}, hold uniformly over all such
triangular arrays.  In particular, the constant in~\eqref{eq:sharp-be} may be chosen
to depend only on $R$.
\end{theorem}

\begin{remark}[Microscopic nonuniversality versus logarithmic universality]
For a fixed nonuniform alphabet distribution, the number of cycles of a
fixed length can converge to a sum of independent geometric variables
rather than to a Poisson variable; this is one of the phenomena identified
by Guerder \cite{Guerder2026}.  Theorem~\ref{thm:standardized} shows that
these nonclassical microscopic laws contribute only $O(1)$ to each fixed
cumulant.  On the expanding logarithmic scale, the universal $1/r$
coefficient dominates.
\end{remark}

\section{Riffle shuffles and major-index bias}

\subsection{Biased and ordinary riffle shuffles}

The inverse description of a generalized riffle shuffle assigns each card
an i.i.d.\ pile label and then sorts stably by that label; see
\cite{BayerDiaconis1992,Stanley2001}.  Depending on convention, this gives
the shuffle permutation or its inverse.  Since inversion preserves cycle
type, this distinction is immaterial here.

\begin{theorem}[One biased riffle shuffle]\label{thm:biased-riffle}
For each $n$, perform one generalized riffle shuffle with pile
probabilities $p^{(n)}=(p_{n,a})$ on a finite or countable ordered set of
piles.  If
\[
 \sup_n\max_a p_{n,a}<1,
\]
then the Gaussian cycle-field theorem, the bivariate functional
Erd\H{o}s--Tur\'an theorem, the endpoint joint law, the sharp
Berry--Esseen estimate~\eqref{eq:sharp-be}, and the Edgeworth expansion
\eqref{eq:edgeworth} all hold uniformly under the stated nonconcentration
bound.
\end{theorem}

\begin{corollary}[One ordinary $a_n$-riffle shuffle]
\label{cor:ordinary-riffle}
Let $a_n\ge2$ be arbitrary, and perform one ordinary $a_n$-riffle shuffle.
Then Theorem~\ref{thm:general-functional} holds.  Moreover, with
$\ell_n=\log n$,
\[
 \Kol\!\left(
 \frac{\log\ord(\rho_n)-\frac12\ell_n^2
       +\ell_n\log\ell_n}
      {\sqrt{\ell_n^3/3}},N(0,1)
 \right)
 \le \frac{C}{\sqrt{\ell_n}},
\]
where $C$ is uniform over all sequences $a_n\ge2$.  With exact
centering by $\Ee\log\ord(\rho_n)$, the Edgeworth expansion
\eqref{eq:edgeworth} and the exact Kolmogorov asymptotic
\eqref{eq:kol-exact} also hold uniformly.
\end{corollary}

\begin{proof}
Use the uniform alphabet distribution on $a_n$ letters.  Its largest atom
is at most $1/2$.
\end{proof}

\begin{proof}[Proof of Theorem~\ref{thm:headline}]
The limit theorem is Corollary~\ref{cor:ordinary-riffle} with $a_n=2$.

For the total-variation statement, the inverse of a one-step $2$-riffle
shuffle has at most one descent.  Let $A_n$ be this event.  Then
\[
 \Pp\{\rho_n^{-1}\in A_n\}=1.
\]
A permutation with at most one descent is determined by choosing the set
of values in its first increasing run, so there are at most $2^n$ such
permutations.  Hence
\[
 \Pp\{\Pi_n\in A_n\}\le\frac{2^n}{n!}\longrightarrow0
\]
for a uniform permutation $\Pi_n$.  Total variation is invariant under
inversion, and therefore tends to one.
\end{proof}

\subsection{Major-index-biased permutations}

For $\pi\in S_n$, write
\[
 \operatorname{maj}(\pi)
 =
 \sum_{j:\pi(j)>\pi(j+1)}j.
\]
Let $\operatorname{Maj}_{n,q}$ be the probability measure proportional to
$q^{\operatorname{maj}(\pi)}$, where $q\in(0,1)$.

\begin{lemma}[Geometric-letter representation]\label{lem:maj-representation}
Let $G_1,\ldots,G_n$ be i.i.d.\ with
\[
 \Pp(G_i=j)=(1-q)q^j,\qquad j\ge0,
\]
and put $\sigma=\operatorname{std}(G_1,\ldots,G_n)$.  If
\[
 w_0=(n,n-1,\ldots,1),
 \qquad
 \tau=w_0\sigma^{-1}w_0,
\]
then $\tau$ has law $\operatorname{Maj}_{n,q}$.  Moreover, $\tau$ and
$\sigma$ have the same cycle type.
\end{lemma}

\begin{proof}
Fix $\sigma\in S_n$, and write $i_k=\sigma^{-1}(k)$.  The event
$\operatorname{std}(G)=\sigma$ is equivalent to
\[
 G_{i_1}\le\cdots\le G_{i_n},
\]
with strict inequality between the $k$th and $(k+1)$st terms whenever
$i_k>i_{k+1}$.  Thus strict inequalities occur at the descent set of
$\sigma^{-1}$.

For a set $D\subseteq[n-1]$, the elementary partition generating function
is
\[
 \sum_{\substack{0\le a_1\le\cdots\le a_n\\
                  a_k<a_{k+1}\ \text{for }k\in D}}
 q^{a_1+\cdots+a_n}
 =
 \frac{q^{\sum_{k\in D}(n-k)}}{\prod_{j=1}^n(1-q^j)}.
\]
It follows that
\[
 \Pp\{\operatorname{std}(G)=\sigma\}
 =
 \frac{q^{\operatorname{comaj}(\sigma^{-1})}}
      {\sum_{\pi\in S_n}q^{\operatorname{comaj}(\pi^{-1})}},
\]
where
\[
 \operatorname{comaj}(\pi)
 =
 \sum_{k\in\operatorname{Des}(\pi)}(n-k).
\]
Conjugation by $w_0$ maps a descent at $k$ to a descent at $n-k$, so
\[
 \operatorname{maj}(w_0\sigma^{-1}w_0)
 =
 \operatorname{comaj}(\sigma^{-1}).
\]
The stated law follows.  Finally, inversion preserves cycle lengths and
conjugation preserves cycle type.
\end{proof}

\begin{corollary}[Major-index Erd\H{o}s--Tur\'an law]
\label{cor:maj}
Let $\pi_{n,q}$ have law $\operatorname{Maj}_{n,q}$ for a fixed
$q\in(0,1)$.  Then Theorems~\ref{thm:cycle-field},
\ref{thm:general-functional}, \ref{thm:sharp-be}, and~\ref{thm:edgeworth},
and Corollary~\ref{cor:endpoint-joint}, hold.  The conclusions, including the
constant in the sharp Berry--Esseen estimate, are uniform for triangular
parameters $q_n\ge q_0>0$.
\end{corollary}

\begin{proof}
The largest atom of the geometric distribution is $1-q$.  Apply
Theorem~\ref{thm:standardized} and
Lemma~\ref{lem:maj-representation}.
\end{proof}

\section{Asymmetric shelf shuffles}

We now turn to a model not generated by ordinary standardization.

An $(m,p)$-shelf shuffle has $m$ labeled shelves.  Cards are dealt
sequentially from the bottom, assigned independently and uniformly to a
shelf, and placed at the top of that shelf with probability $p$ and at the
bottom with probability $q=1-p$.  The shelf piles are then concatenated in
shelf order.  Let $\nu_{n,m}^{(p)}$ be the induced law on $S_n$.

Tripathi proved the following exact cycle-index formula.

\begin{inputtheorem}[Tripathi's shelf-shuffle cycle index
\cite{Tripathi2026}]\label{input:tripathi}
For fixed $m\ge1$ and $p\in(0,1)$,
\[
 1+\sum_{n\ge1}u^n
 \sum_{\pi\in S_n}\nu_{n,m}^{(p)}(\pi)
 \prod_{i\ge1}x_i^{C_i(\pi)}
 =
 \exp\left\{
 \sum_{i\ge1}\sum_{r\ge1}
 \frac{B_{i,r}^{(m,p)}}{ir}x_i^ru^{ir}
 \right\},
\]
where
\begin{equation}\label{eq:tripathi-B}
 B_{i,r}^{(m,p)}
 =
 \sum_{d\mid i}
 \mu(d)m^{i/d-ir}
 \bigl(p^{dr}-(-q)^{dr}\bigr)^{i/d}.
\end{equation}
At $\mathbf x=\mathbf1$, the generating function equals $(1-u)^{-1}$.
\end{inputtheorem}

\begin{lemma}[Uniform exponential stability for shelf shuffles]
\label{lem:shelf-stability}
Fix $\eps\in(0,1/2]$.  There are $A_\eps<\infty$ and
$\rho_\eps\in(0,1)$ such that, for all
\[
 m\ge1,\qquad \eps\le p\le1-\eps,
\]
the coefficients in \eqref{eq:tripathi-B} satisfy
\[
 |B_{i,1}^{(m,p)}-1|\le A_\eps\rho_\eps^i,
\]
and
\[
 |B_{i,r}^{(m,p)}|\le A_\eps\rho_\eps^{ir},
 \qquad r\ge2.
\]
\end{lemma}

\begin{proof}
Put $q=1-p$ and
\[
 \alpha(p)=\sqrt{p^2+q^2}.
\]Uniformly for $p\in[\eps,1-\eps]$,
\[
 \alpha(p)\le
 \alpha_\eps:=\sqrt{1-2\eps(1-\eps)}<1.
\]
For every integer $s\ge2$, monotonicity of $\ell^r$ norms gives
\[
 p^s+q^s\le\alpha(p)^s.
\]

When $r=1$ and $d=1$, the summand in
\eqref{eq:tripathi-B} is
\[
 (p+q)^i=1.
\]
For $r=1$ and $d\ge2$,
\[
 \left|
 m^{i/d-i}\bigl(p^d-(-q)^d\bigr)^{i/d}
 \right|
 \le\alpha_\eps^i.
\]
For $r\ge2$ and every $d\mid i$,
\[
 \left|
 m^{i/d-ir}\bigl(p^{dr}-(-q)^{dr}\bigr)^{i/d}
 \right|
 \le\alpha_\eps^{ir},
\]
because $m^{1-dr}\le1$.  There are at most $i$ divisors.  Choose
$\rho_\eps\in(\alpha_\eps,1)$ and absorb the divisor factor into
$A_\eps\rho_\eps^{ir}$.
\end{proof}

\begin{theorem}[One asymmetric shelf shuffle]
\label{thm:shelf}
Let
\[
 m_n\ge1,
 \qquad
 \eps\le p_n\le1-\eps,
\]
and let $\pi_n$ have law $\nu_{n,m_n}^{(p_n)}$.  Then the Gaussian
cycle-field theorem, the bivariate functional Erd\H{o}s--Tur\'an theorem,
the endpoint joint law, the sharp Berry--Esseen theorem, and the universal
first Edgeworth expansion hold uniformly
over these parameters.  In particular, with $\ell_n=\log n$,
\[
 \Kol\!\left(
 \frac{\log\ord(\pi_n)-\frac12\ell_n^2
       +\ell_n\log\ell_n}
      {\sqrt{\ell_n^3/3}},N(0,1)
 \right)
 \le \frac{C_\eps}{\sqrt{\ell_n}}.
\]
\end{theorem}

\begin{proof}
For each target $n$, use the cycle-index family with the same parameters
$(m_n,p_n)$ at every auxiliary size.  Input
Theorem~\ref{input:tripathi} and
Lemma~\ref{lem:shelf-stability} verify exponential Ewens stability.
Apply Theorems~\ref{thm:cycle-field},
\ref{thm:general-functional}, \ref{thm:sharp-be}, and~\ref{thm:edgeworth}.
\end{proof}

\begin{remark}
The theorem is uniform in the number of shelves.  It includes a single
shelf, a fixed number of shelves, and an arbitrarily growing number of
shelves.  The parameter restriction excludes only degeneration toward
always placing at the top or always placing at the bottom.
\end{remark}

\section{Boundary of the universality class}

The nonconcentration conditions cannot simply be removed.

For standardized permutations, take a two-letter alphabet with
\[
 p_{n,0}=1-\delta_n,\qquad p_{n,1}=\delta_n,
\]
where $n\delta_n\to0$.  With probability tending to one, every letter is
zero, so the standardized permutation is the identity and its order is
one.  Thus $\|p^{(n)}\|_\infty\to1$ can destroy the Erd\H{o}s--Tur\'an
law completely.

For shelf shuffles, the endpoint parameters can also destroy the
conclusion.  When $m=1$, the rules $p=0$ and $p=1$ produce, respectively,
the identity and the reversal permutation (up to the harmless convention
for reading the deck); their orders are bounded by two.  Thus the theorem
identifies a genuine uniform nondegeneracy regime.  It does not assert
universality throughout every possible boundary scaling.

A natural next question is to classify the transition when the exponential
stability parameter $\rho=\rho_n$ tends to one.  In the standardized model
this means that one letter becomes dominant.  In the shelf model it means
$p_n$ approaches zero or one.  Such regimes may exhibit different order
scales and are not addressed here.

\section{Consequences and further directions}

The results establish a reusable principle:

\begin{quote}
An exponentially localized perturbation of the Ewens cycle index may
change every fixed short-cycle law while leaving the entire logarithmic
Gaussian cycle field, and hence the Erd\H{o}s--Tur\'an order process,
unchanged.
\end{quote}

The theorem applies with no convergence of the underlying parameter
profile.  It is therefore robust under triangular arrays of biased riffle
probabilities, alphabet sizes, shelf numbers, and shelf asymmetries within
the stated compact regime.

The scalar order-level problem is settled through the first Edgeworth term:
the sharp rate, its optimality, and the universal leading Kolmogorov
constant are all identified.  The long-cycle total-variation theorem also
suggests quantitative finite-dimensional extensions, but a sharp
multivariate or path-space Berry--Esseen theorem requires a choice of metric
and a separate Gaussian approximation argument; it is not claimed here.
A second direction is to weaken exponential stability to polynomial or
operator-theoretic mixing of the cycle-index coefficients.  That extension
would be relevant to dense kernel/permanental permutations and weak
Mallows permutations.

\appendix

\section{A coefficient-tail lemma}

We record the elementary estimate used in
Proposition~\ref{prop:mixed-cumulant}.

\begin{lemma}\label{lem:coefficient-tail}
Let $q$ be fixed.  Suppose
\[
 A_j(u)=\sum_{r\le U}a_{j,r}u^r,
 \qquad
 \sum_r|a_{j,r}|\le M\log n,
\]
and let
\[
 E_j(u)=\sum_{m\ge b}e_{j,m}u^m,
 \qquad
 \sum_{m\ge b}|e_{j,m}|\le M\rho^b
\]
for $1\le j\le q$, where $\rho<1$.  If $qU<n$, then the difference between
the sum of coefficients through degree $n$ and the value at $u=1$ of any
product of at most $q$ factors selected from the $A_j+E_j$ is
\[
 O_{q,M}\bigl((1+\log n)^q\rho^b\bigr).
\]
\end{lemma}

\begin{proof}
The product containing only $A$ factors has degree at most $qU<n$ and
therefore contributes no tail.  Every other product contains an $E$
factor.  Bound its total absolute coefficient mass by the product of the
$\ell^1$ coefficient norms.
\end{proof}

\section{A tightness criterion with vanishing errors}

The following standard variant of Billingsley's two-increment criterion is
convenient for triangular arrays.

\begin{lemma}\label{lem:tightness-criterion}
Let $X_n$ be cadlag processes.  Suppose their finite-dimensional
distributions are tight, their maximal jumps converge to zero in
probability, and there are $C<\infty$ and $\eta_n\to0$ such that
\[
 \Ee\left[
 |X_n(t)-X_n(s)|^2
 |X_n(u)-X_n(t)|^2
 \right]
 \le C(u-s)^2+\eta_n
\]
for all $0\le s\le t\le u\le1$.  Then $(X_n)$ is tight in $D[0,1]$, and
every subsequential limit is continuous.
\end{lemma}

\begin{proof}
Fix a mesh width $\delta>0$.  In the proof of Billingsley's
two-increment criterion, the interval $[0,1]$ is first covered by a finite
mesh depending only on $\delta$, and Markov's inequality is applied to
pairs of adjacent increments.  Replacing the usual bound
$C(u-s)^2$ by $C(u-s)^2+\eta_n$ adds at most
$C_\delta\eta_n$ to the resulting probability estimate, where
$C_\delta<\infty$ is independent of $n$.  This additional term vanishes
when $n\to\infty$ with $\delta$ fixed.  The remainder of the proof of
\cite[Theorem 13.5]{Billingsley1999} is unchanged and gives the required
cadlag modulus after $\delta\downarrow0$.  The maximal-jump assumption
rules out discontinuities in every subsequential limit.
\end{proof}

\begin{thebibliography}{99}

\bibitem{ABT2003}
R. Arratia, A. D. Barbour, and S. Tavar\'e,
\newblock \emph{Logarithmic Combinatorial Structures: A Probabilistic
Approach},
\newblock EMS Monographs in Mathematics, European Mathematical Society,
Z\"urich, 2003.
\newblock \href{https://doi.org/10.4171/000}{doi:10.4171/000}.

\bibitem{ArratiaTavare1992}
R. Arratia and S. Tavar\'e,
\newblock Limit theorems for combinatorial structures via discrete process
approximations,
\newblock \emph{Random Structures \& Algorithms} 3 (1992), no.~3, 321--345.
\newblock \href{https://doi.org/10.1002/rsa.3240030310}
{doi:10.1002/rsa.3240030310}.

\bibitem{BarbourTavare1994}
A. D. Barbour and S. Tavar\'e,
\newblock A rate for the Erd\H{o}s--Tur\'an law,
\newblock \emph{Combin. Probab. Comput.} 3 (1994), no.~2, 167--176.
\newblock \href{https://doi.org/10.1017/S0963548300001097}
{doi:10.1017/S0963548300001097}.

\bibitem{BayerDiaconis1992}
D. Bayer and P. Diaconis,
\newblock Trailing the dovetail shuffle to its lair,
\newblock \emph{Ann. Appl. Probab.} 2 (1992), no.~2, 294--313.
\newblock \href{https://doi.org/10.1214/aoap/1177005705}
{doi:10.1214/aoap/1177005705}.

\bibitem{Billingsley1999}
P. Billingsley,
\newblock \emph{Convergence of Probability Measures},
\newblock second edition, Wiley, New York, 1999.

\bibitem{DiaconisFulmanHolmes2013}
P. Diaconis, J. Fulman, and S. Holmes,
\newblock Analysis of casino shelf shuffling machines,
\newblock \emph{Ann. Appl. Probab.} 23 (2013), no.~4, 1692--1720.
\newblock \href{https://doi.org/10.1214/12-AAP884}
{doi:10.1214/12-AAP884}.

\bibitem{DiaconisMcGrathPitman1995}
P. Diaconis, M. McGrath, and J. Pitman,
\newblock Riffle shuffles, cycles, and descents,
\newblock \emph{Combinatorica} 15 (1995), no.~1, 11--29.
\newblock \href{https://doi.org/10.1007/BF01294457}
{doi:10.1007/BF01294457}.

\bibitem{ErdosTuran1967}
P. Erd\H{o}s and P. Tur\'an,
\newblock On some problems of a statistical group theory. III,
\newblock \emph{Acta Math. Acad. Sci. Hungar.} 18 (1967), 309--320.

\bibitem{Guerder2026}
A. Guerder,
\newblock Cycle structure of random standardized permutations,
\newblock in \emph{37th International Conference on Probabilistic,
Combinatorial and Asymptotic Methods for the Analysis of Algorithms
(AofA 2026)}, LIPIcs 381 (2026), Article 15, 1--16.
\newblock \href{https://doi.org/10.4230/LIPIcs.AofA.2026.15}
{doi:10.4230/LIPIcs.AofA.2026.15};
full version \href{https://arxiv.org/abs/2603.24127}{arXiv:2603.24127}.

\bibitem{HardyWright2008}
G. H. Hardy and E. M. Wright,
\newblock \emph{An Introduction to the Theory of Numbers},
\newblock sixth edition, Oxford University Press, Oxford, 2008.

\bibitem{Stanley2001}
R. P. Stanley,
\newblock Generalized riffle shuffles and quasisymmetric functions,
\newblock \emph{Ann. Comb.} 5 (2001), no.~3--4, 479--491.
\newblock \href{https://doi.org/10.1007/s00026-001-8023-7}
{doi:10.1007/s00026-001-8023-7}.

\bibitem{Tripathi2026}
R. Tripathi,
\newblock Analysis of the asymmetric shelf shuffle,
\newblock preprint, 2026.
\newblock \href{https://arxiv.org/abs/2606.18047}{arXiv:2606.18047}.

\bibitem{Zacharovas2009}
V. Zacharovas,
\newblock \emph{Distribution of Random Variables on the Symmetric Group},
\newblock doctoral dissertation, Vilnius University, 2004;
\newblock \href{https://arxiv.org/abs/0901.1733}{arXiv:0901.1733}.

\end{thebibliography}

\end{document}